La validez de los resultados del enfoque investigador del proyecto se basa en la comparabilidad de estos mismos, y , por ello, la validación de estos debería de responder a las siguientes preguntas:

-¿Qué problema había antes?

-¿Qué se hizo para validar la evaluación?

-¿Por qué esto hace que los resultados sean comparables?

-¿Qué se corrigió respecto a versiones anteriores?

-¿Qué papel tiene el baseline?


\textbf{Protocolo de evaluación utilizado}:

El objetivo de crear un protocolo de evaluación no ha sido solo tener un proceso de ejecución sin errores, si no que también, y realmente de más importancia, asegurar que las métricas obtenidas reflejasen realmente las capacidades de cada función de pérdida en el modelo.
El protocolo de evaluación propuesto es el siguiente:

\begin{itemize}
    \item 1. Uso de \textit{dataset} / \textit{representative\_test}.
    \item 2. Evaluación de todos los modelos sobre las mismas pistas.
    \item 3. Normalización de referencias y estimaciones a mono-canal.
    \item 4. Se recortan las pistas a una longitud común
    \item 5. Se usa la misma función de despliegue del modelo que en \textit{deployment}.
    \item 6. Se guarda un sumario en formato \textit{JSON} para cada modelo.
\end{itemize}

Esto se ha hecho siguiendo las recomendaciones habituales de reproducibilidad en \textit{Machine Learning} \cite{MachineLearningPaperReproducibilityChecklist}:
documentar dataset, \textit{splits}, preprocesado, hiperparámetros y resultados. La checklist de reproducibilidad incluye precisamente detalles del dataset,
particiones \textit{train}/\textit{validation}/\textit{test}, datos excluidos y pasos de preprocesado.

Inicialmente el proceso, aunque generase archivos, esto no significaba que los resultados fuesen realmente los buscados, en particular, porque las pistas de salida eran exactamente iguales a las pistas de entrada en la mayoría de los casos, por lo que algo se estaba haciendo mal.
Para dar alguna medida de validez a los resultados obtenidos, se trasladó el conjunto de test a un test fijo, se creó la validación de \textit{baseline}, se añadió una sección en el código donde se normalizaban las pistas en mono-canal, recorte de pistas a una longitud común, impresión de métricas de análisis de las máscaras por pantalla en el entrenamiento y normalización del flujo de inferencia para todas las funciones de pérdida, donde previamente existían variaciones.
A partir de aquí, todos los modelos pasan a evaluarse sobre las mismas pistas, con el mismo preprocesado, la misma métrica de evaluación y el mismo flujo de inferencia, además , añadir la métrica \textit{baseline} ayuda a identificar los resultados que generaban audios pero no separaban mas allá de copiar la mezcla de entrada.

El protocolo de evaluación final queda definido así:

\begin{enumerate}
    \item Cargar el conjunto de test fijo \textit{representative\_test}.
    \item Evaluar todos los modelos sobre las mismas pistas. 
    \item Ejecutar la inferencia mediante la misma función utilizada en el despliegue. 
    \item Convertir referencias y estimaciones a mono-canal. 
    \item Recortar referencias y estimaciones a una longitud común. 
    \item Mantener un orden fijo de fuentes para dos y cuatro \textit{stems}. 
    \item Calcular SI-SDR para cada modelo y pista. 
    \item Comparar los resultados frente al baseline de mezcla. 
    \item Guardar los resultados y sumarios de evaluación en archivos estructurados. 
\end{enumerate}